\documentclass[12pt]{article}

\usepackage{graphicx}

\usepackage{makecell}
\usepackage[figuresright]{rotating}
\usepackage{multicol, caption}
\usepackage{authblk}
\usepackage[margin=1.0in]{geometry}
\usepackage[export]{adjustbox}

\textheight=210mm
\textwidth=155mm

\hoffset=-1mm
\voffset=-10mm

\begin{document}
Submission MS \#CHA19-AR-00322


\begin{center}
{\bf Response to comments of reviewer \#1}\\
\vspace{2mm}
Concerning paper ``Generalised early warning
signals in multivariate and gridded data with an application to
tropical cyclones''\\[5pt]
by {\it J.~Prettyman, T.~Kuna, V.~N.~Livina}
\end{center}

\vspace{3mm}

We thank the reviewer for the time taken to read through this submission and provide positive constructive feedback with suggestions for improvement. We will reply to the comments here in the order in which they are given.

\begin{enumerate}
\item \textit{...there is still a lack of explanations in the main text regarding the tropical cyclone examples. I suggest to emphasize explicitly in the paper that this example has nothing to do with the bifurcation case described in the previous sections. It should be stated clearly in motivation part that it was shown in some works (e.g. [Ashwin et al.]) that the proposed approach can be useful even for non-bifurcation tipping points, and here an attempt is made to apply it to data where such a case is possible. In particular, the replies of the authors to my comments 1 and 2 could be used in the paper text.}

\textbf{reply (1):} Thank you for this suggestion. It may have been an oversight on our part that we devoted some time defending our choice in our previous reply but did not think to include this in the manuscript. We certainly agree that this is a useful addition and have included the argument now (bottom of page 4), including several references. 

\item \textit{Next, a relation of the considered situations to Ashwin's et al. cases should be shown and mathematical reasons why the methods can work in the tropical cyclone case (e.g. for the model proposed in the paper) should be outlined. I think, that the relevancy of the cyclone example to the method used remains unclear without such explanations. }
	
\textbf{reply (1):} The argument on pages 4/5 also includes a reference to the general applicability of the critical slowing down phenomenon and, therefore, autocorrelation-based techniques in a range of dynamical systems and tipping scenarios. We have noted in the same paragraph that the chosen system appears to exhibit a transitional rather than bifurcational tipping point. This is also now mentioned on page 7 when referring to the model.
	
\item \textit{I would like to clarify what I meant in comment 9 (I do not insist upon changing the authors' derivation in the paper). We can multiply the left and right sides of Eq. 2 by their transposes and average them over n. We get an equation for the covariance matrix $C=BCB^T+SS^T$. To solve this we can iterate $C_{n+1}=BC_nB^T+SS^T$. It is easy to see that at infinite n it converges right to the Eq. 6 independently on initial $C_0$.}

\textbf{reply (3):}	Thank you for this clarification. We have chosen not to change the derivation in the paper since it has already been written and remains valid. However, it is useful to have seen this alternative approach and we are grateful that reviewer has taken the time to so thoroughly analyse this work.

\end{enumerate}

We thank the reviewer again for their time in responding carefully to this submission and the previous iterations. As a result of this conversation we feel that the strength of the paper, and also our own understanding of the field, has improved considerably. We hope that the revised manuscript is now suitable for publication in the Chaos magazine and that research in this field continues to expand.

\vspace{5mm}

\noindent
Yours faithfully,

\vspace{1mm}

\noindent
\it J.~Prettyman, T.~Kuna, V.~N.~Livina

\end{document}
